% Unified matrix-form reformulation of the current OODKA-OT implementation.
% The source of truth for content is docs/oodka_current_method.tex; this file
% reorganizes the same executable method around a small number of master
% operators and matrix equations.

\documentclass[10pt]{article}

\usepackage[margin=1in]{geometry}
\usepackage[T1]{fontenc}
\usepackage[utf8]{inputenc}
\usepackage{amsmath,amssymb,mathtools}
\usepackage{bm}
\usepackage{booktabs}
\usepackage{graphicx}
\usepackage{microtype}
\usepackage[hidelinks]{hyperref}

\newcommand{\R}{\mathbb{R}}
\newcommand{\cI}{\mathcal{I}}
\newcommand{\cQ}{\mathcal{Q}}
\newcommand{\cL}{\mathcal{L}}
\newcommand{\cG}{\mathcal{G}}
\newcommand{\eps}{\varepsilon}
\newcommand{\epsnum}{\varepsilon_{\mathrm{num}}}
\newcommand{\ones}{\bm{1}}
\newcommand{\diag}{\operatorname{Diag}}
\newcommand{\Mat}{\operatorname{Mat}}
\newcommand{\Resize}{\operatorname{Resize}}
\newcommand{\IN}{\operatorname{IN}}
\newcommand{\LN}{\operatorname{LN}}
\newcommand{\GELU}{\operatorname{GELU}}
\newcommand{\softplus}{\operatorname{softplus}}
\newcommand{\sigm}{\operatorname{sigmoid}}
\newcommand{\Norm}{\operatorname{Norm}}
\newcommand{\clip}{\operatorname{clip}}
\newcommand{\KL}{\operatorname{KL}}
\newcommand{\stopgrad}{\operatorname{stopgrad}}
\newcommand{\rowcos}{\operatorname{rowcos}}
\newcommand{\Ent}{\operatorname{Ent}}
\newcommand{\RoiBox}{\operatorname{Box}}
\newcommand{\Crop}{\operatorname{CropResize}}
\newcommand{\relu}[1]{\left[#1\right]_{+}}
\newcommand{\ind}{\mathbb{1}}
\newcommand{\balpha}{\bm{\alpha}}
\newcommand{\bbeta}{\bm{\beta}}

\title{OODKA-OT: A Unified Matrix Formulation}
\date{}

\begin{document}
\maketitle

\begin{abstract}
We present the current OODKA-OT implementation as one coupled matrix system.
A frozen BiomedParse student representation and a training-only frozen nnUNet
expert are factorized into a structure-oriented component $P$ and a
task-residual component $S$.  The two feature factors are constrained by masked
reconstruction and cross-correlation suppression, aligned by one branch-indexed
optimal-transport problem, and recombined by a prompt-conditioned stochastic
matrix field.  The same student operator is composed twice for LGE inference:
first globally for anatomy localization and then locally inside a predicted
myocardial region.  This presentation changes no executable behavior; it only
collects the implementation into a smaller set of reusable mathematical
operators.
\end{abstract}

\section{A Single Mathematical View of the System}

Let $B_{\phi_b}$ denote the frozen BiomedParse model and $E_{\phi_n}$ the
frozen nnUNet expert.  Only the feature factorization modules and the spatial
router are trainable; their parameters are collected in $\Theta$, while
$\phi_b$ and $\phi_n$ are fixed.  The expert and all transport operators exist
only during training.

We keep the following indices fixed throughout the presentation:
\begin{center}
\begin{tabular}{cl}
\toprule
Index & Meaning \\
\midrule
$d\in\{b,n\}$ & foundation (BiomedParse) or expert (nnUNet) side \\
$q,\xi\in\{p,s\}$ & structure-oriented or task-residual factor \\
$\ell\in\cI$ & encoder feature level \\
$r\in\cQ$ & frozen pixel-decoder output resolution \\
$j\in\{A,R\}$ & full-image anatomy or ROI-refinement pass \\
$c$ & prompt/class index \\
\bottomrule
\end{tabular}
\end{center}

Training uses $B$ independent blocks with $Z$ slices per block.  At the four
implemented feature levels
\begin{equation}
    \cI=\{2,3,4,5\},
\end{equation}
the frozen encoders return foundation features $\bm Z^b_\ell$ and raw expert
features $\bm Z^{n,\mathrm{raw}}_\ell$.  The latter are trilinearly aligned to the
foundation lattice,
\begin{equation}
    \bm Z^n_\ell
    =\Resize_{\mathrm{tri}}
      (\bm Z^{n,\mathrm{raw}}_\ell;Z,H_\ell,W_\ell).
\end{equation}
We unfold a five-dimensional tensor into a row-token matrix by placing channels
in columns:
\begin{equation}
    \bm Z^{d,\mathrm{mat}}_\ell=\Mat(\bm Z^d_\ell)
       \in\R^{N_\ell\times C^d_\ell},
    \qquad N_\ell=BZH_\ell W_\ell,
    \qquad d\in\{b,n\}.
    \label{eq:matrix-features}
\end{equation}
Although written as one matrix, $\Mat$ retains the partition into $B$ blocks.
Every instance-normalization symbol below denotes the induced block operator
\begin{equation}
    \IN^{\mathrm{blk}}
    =\Mat\circ\operatorname{InstanceNorm3d}\circ\Mat^{-1},
\end{equation}
which normalizes each block and channel over its own $ZHW$ sites, never across
different blocks.  Repeated tail sites still participate in these forward
normalization statistics, exactly as in the fixed-shape implementation.

The row ordering is $(b,z,\bm u)$.  A diagonal matrix
$\bm V_\ell\in\{0,1\}^{N_\ell\times N_\ell}$ repeats the valid-slice flag
$v_{bz}$ over all spatial sites.  It therefore removes repeated tail slices
without changing tensor shapes.  In the LGE path $Z=1$ and $\bm V_\ell$ is
normally the identity.  Invalid tail targets are stored as $-1$ and are
excluded from segmentation, reconstruction, separation, and OT, although the
fixed-shape forward still evaluates their repeated inputs.  We use $\bm V_r$
for the identical broadcast construction on a decoder-output grid $r$.
When a block-local matrix is required, $\bm J_{\ell b}$ selects the rows of
block $b$:
\begin{equation}
 \bm P^d_{\ell b}=\bm J_{\ell b}\bm P^d_\ell,\qquad
 \bm V_{\ell b}=\bm J_{\ell b}\bm V_\ell\bm J_{\ell b}^{\top},
\end{equation}
and analogously for $\bm S^d_{\ell b}$ and other factor matrices.

For generic 2.5D training, the three BiomedParse channels associated with
center slice $z$ are $(x_{z-1},x_z,x_{z+1})$, with boundary replication.  LGE
uses $Z=1$ and the pseudo-RGB triplet $(x_z,x_z,x_z)$.  The generic CT path
uses level $40$ and width $400$, equivalently
\begin{equation}
    x_i^b=255\,[\clip(x_i,-160,240)+160]/400,
\end{equation}
before the frozen BiomedParse pixel normalization.

The training computation can be summarized before introducing its components as
\begin{equation}
\boxed{
\begin{aligned}
 \mathsf{Factorization}:\quad
 &(\bm P^b,\bm S^b;\bm P^n,\bm S^n)
 =\mathcal D_\Theta(\bm Z^{b,\mathrm{mat}},\bm Z^{n,\mathrm{mat}}),\\
 \mathsf{Prediction}:\quad
 &(\bm P^b,\bm S^b,\bm T)
 \xrightarrow{\ \bm\Gamma_\Theta\ }
 \widetilde{\bm F}
 \xrightarrow{\ \mathcal H_{\phi_b}\ }\bm L_\Theta,\\
 \mathsf{Supervision}:\quad
 &(\bm P^b,\bm S^b;\bm P^n,\bm S^n)
 \xrightarrow{\ P\text{-OT}/S\text{-UOT}\ }
 (\widetilde{\bm Y}^{p},\widetilde{\bm Y}^{s}).
\end{aligned}}
\label{eq:master-flow}
\end{equation}
Here and below, bold tuples without a level subscript collect all
$\ell\in\cI$.

The deployed student can already be stated independently of its training-only
teacher.  Let $\bm T$ be the matrix of prompt embeddings,
$\bm F_q=\Psi_{\phi_b}(\mathcal D^b_{\Theta,q}(B_{\phi_b}(\bm x^b)))$ the
frozen pixel-decoder pyramid induced by branch $q\in\{p,s\}$,
$\mathcal M_{\bm\Gamma}$ the spatial mixing operator, $\mathcal H_{\phi_b}$
the frozen prompt predictor, $\mathcal A_{Q_m}$ query averaging, and
$\mathcal U$ output-logit resizing.  The entire inference map is
\begin{equation}
\boxed{
 \bm L_\Theta(\bm x^b,\bm T)
 =\mathcal U\circ\mathcal A_{Q_m}\circ\mathcal H_{\phi_b}
 \left(
   \mathcal M_{\bm\Gamma_\Theta(\bm T)}[\bm F_p,\bm F_s],\bm T
 \right).}
\label{eq:master-student}
\end{equation}
The subsequent terms either constrain the deployed factors in
Eq.~\eqref{eq:master-student} or construct their training-only teacher; the
nnUNet expert is not an argument of this deployed map.

\section{Multi-scale Factorization into Complementary Feature Components}

The pointwise $1\times1\times1$ convolutions are channel matrices after the
unfolding in Eq.~\eqref{eq:matrix-features}.  Consequently, the foundation-side
factorization at level $\ell$ is the single block-matrix equation
\begin{equation}
    \boxed{
    \begin{bmatrix}\bm P^b_\ell&\bm S^b_\ell\end{bmatrix}
    =\bm Z^{b,\mathrm{mat}}_\ell
     \begin{bmatrix}\bm W^b_{\ell p}&\bm W^b_{\ell s}\end{bmatrix}.}
    \label{eq:foundation-block-factorization}
\end{equation}
Both channel matrices are bias-free, preserve $C^b_\ell$ channels, and are
independent:
$\bm W^b_{\ell p},\bm W^b_{\ell s}\in
\R^{C^b_\ell\times C^b_\ell}$.  No activation or normalization follows either
projection.

The expert first enters the same channel space through an intermediate matrix
\begin{equation}
    \bm U_\ell
    =\GELU\!\left(\IN_\ell(\bm Z^{n,\mathrm{mat}}_\ell\bm W^0_\ell)\right),
\end{equation}
and is then factorized as
\begin{equation}
    \boxed{
    \begin{bmatrix}\bm P^n_\ell&\bm S^n_\ell\end{bmatrix}
    =\begin{bmatrix}
       \mathcal N_{\ell p}(\bm U_\ell\bm W^n_{\ell p})&
       \mathcal N_{\ell s}(\bm U_\ell\bm W^n_{\ell s})
     \end{bmatrix}.}
    \label{eq:expert-block-factorization}
\end{equation}
Here $\bm W^0_\ell\in\R^{C^n_\ell\times C^{\mathrm{mid}}_\ell}$ and
$\bm W^n_{\ell p},\bm W^n_{\ell s}
\in\R^{C^{\mathrm{mid}}_\ell\times C^b_\ell}$ align the expert branches to
the foundation channel dimension.
The operators $\mathcal N_{\ell p}$ and $\mathcal N_{\ell s}$ are affine
instance normalization for $\ell=2,3,4$ and identities for $\ell=5$ in the
default configuration.  Equations~\eqref{eq:foundation-block-factorization}
and~\eqref{eq:expert-block-factorization} do not assign semantics by
architecture alone: the $P$ branch becomes structure-oriented through balanced
OT, whereas the $S$ branch becomes task-residual through error-aware UOT.

\subsection{One masked relative-energy operator}

All reconstruction terms use the same masked relative Frobenius energy.  For
matrices with $C$ columns, define
\begin{equation}
    \mathfrak R(\bm A,\bm X^\star;\bm V)
    =\frac{\|\bm V(\bm A-\bm X^\star)\|_F^2/n_V}
           {\|\bm V\bm X^\star\|_F^2/n_V+10^{-8}},
    \qquad n_V=CHW\max\!\left(\sum_{b,z}v_{bz},1\right).
    \label{eq:masked-relative-energy}
\end{equation}
This is exactly the valid-slice energy-normalized MSE, expressed without a
per-element summation.  Let $\mathcal R_{\ell p}$ and
$\mathcal R_{\ell s}$ be the expert decoders
Conv--IN--GELU--Conv.  The two level-wise reconstructions are
\begin{equation}
\begin{aligned}
    \widehat{\bm Z}^{b,\mathrm{mat}}_\ell&=\bm P^b_\ell+\bm S^b_\ell,\\
    \widehat{\bm Z}^{n,\mathrm{mat}}_\ell
      &=\mathcal R_{\ell p}(\bm P^n_\ell)
        +\mathcal R_{\ell s}(\bm S^n_\ell).
\end{aligned}
\label{eq:matrix-reconstructions}
\end{equation}

Let $\Psi^r_{\phi_b}(\bm P)$ denote the frozen BiomedParse pixel-decoder output
at $r\in\cQ=\{\mathrm{mask},5,4,3\}$ when the complete multi-scale tuple
$\bm P=\{\bm P_\ell\}_{\ell\in\cI}$ is injected.  Set
\begin{equation}
    \bm F_p^r=\Psi^r_{\phi_b}(\bm P^b),\qquad
    \bm F_s^r=\Psi^r_{\phi_b}(\bm S^b),\qquad
    \bm F_{p+s}^r=\Psi^r_{\phi_b}(\bm P^b+\bm S^b).
\end{equation}
The complete autoencoding objective is then one operator average,
\begin{equation}
\boxed{
\cL_{\mathrm{ae}}
=\frac{1}{9}\left\{
 \sum_{\ell\in\cI}\sum_{d\in\{b,n\}}
 \mathfrak R(\widehat{\bm Z}^{d,\mathrm{mat}}_\ell,
              \bm Z^{d,\mathrm{mat}}_\ell;\bm V_\ell)
 +\frac{1}{|\cQ|}\sum_{r\in\cQ}
 \mathfrak R(\bm F_p^r+\bm F_s^r,\bm F_{p+s}^r;\bm V_r)
 \right\}.}
\label{eq:unified-reconstruction}
\end{equation}
Thus the eight encoder-level terms and one already-averaged decoder-consistency
term form the nine implementation terms.

\subsection{Separation as a cross-correlation matrix norm}

Let $\mathcal Z_{\bm V_b}(\bm X)$ denote block-wise channel standardization over
valid spatial rows: masked mean centering followed by channel standardization,
with invalid rows set to zero.  It follows the code's sample-standard-deviation
path when every slice is valid and its masked population-RMS path otherwise.
For model side $d\in\{b,n\}$, define the cross-correlation matrix
\begin{equation}
    \bm C^{d}_{\ell b}
    =\frac{1}{N_{\ell b}^{\mathrm{valid}}}
      \mathcal Z_{\bm V_{\ell b}}(\bm P^d_{\ell b})^{\!\top}
      \mathcal Z_{\bm V_{\ell b}}(\bm S^d_{\ell b})
      \in\R^{C^b_\ell\times C^b_\ell}.
\end{equation}
The implemented ``orthogonality'' loss is precisely the entrywise matrix
$\ell_1$ norm
\begin{equation}
    \boxed{
    \cL_{\mathrm{ort}}
    =\sum_{\ell\in\cI}\sum_{d\in\{b,n\}}
      \frac{1}{B(C^b_\ell)^2}
      \sum_{b=1}^{B}\|\bm C^{d}_{\ell b}\|_{1,1}.}
    \label{eq:matrix-separation}
\end{equation}
It suppresses every $P$--$S$ channel pair, not merely corresponding channels.
It is a second-order decorrelation penalty rather than a claim of statistical
independence; reconstruction, segmentation, and transport prevent the trivial
zero-feature solution and determine the two branches' operational roles.

\section{Prompt-conditioned Beta Routing as a Stochastic Matrix Field}

Stack the frozen class embeddings into
$\bm T=[\bm t_1^\top;\ldots;\bm t_{C_{\mathrm{pr}}}^\top]
\in\R^{C_{\mathrm{pr}}\times d_t}$.  The router computes all prompts at once,
\begin{equation}
    \bm H=\GELU\!\left(\LN(\bm T)\bm W_h
       +\ones\bm b_h^\top\right),
    \qquad
    \bm A^\alpha=\operatorname{Head}_\alpha(\bm H),\quad
    \bm A^\beta=\operatorname{Head}_\beta(\bm H),
    \label{eq:router-coefficients}
\end{equation}
where $\bm A^\alpha,\bm A^\beta\in\R^{C_{\mathrm{pr}}\times K}$.

On the mask-feature grid with $N_0=N_{\mathrm{mask}}$ sites, collect the
$K=G^2$ Gaussian radial bases in the design matrix
$\bm\Phi_0\in\R^{N_0\times K}$,
\begin{equation}
    [\bm\Phi_0]_{ik}
    =\exp\!\left[-\frac{\|\bm u_i-\bm\mu_k\|_2^2}
                         {2\sigma_{\mathrm{rbf}}^2}\right],
    \quad \bm u_i,\bm\mu_k\in[-1,1]^2.
\end{equation}
The current defaults are $G=8$, $K=64$, and
$\sigma_{\mathrm{rbf}}=3/(G-1)$.  Instead of predicting pixels separately,
the router constructs both Beta-parameter fields by two matrix products:
\begin{equation}
\boxed{
\begin{aligned}
 \bm\eta^\alpha
   &=b_\alpha\ones_{C_{\mathrm{pr}}}\ones_{N_0}^\top
     +K^{-1/2}\bm A^\alpha\bm\Phi_0^\top,
 &\balpha&=\ones_{C_{\mathrm{pr}}}\ones_{N_0}^\top
             +\softplus(\bm\eta^\alpha),\\
 \bm\eta^\beta
   &=b_\beta\ones_{C_{\mathrm{pr}}}\ones_{N_0}^\top
     +K^{-1/2}\bm A^\beta\bm\Phi_0^\top,
 &\bbeta&=\ones_{C_{\mathrm{pr}}}\ones_{N_0}^\top
            +\softplus(\bm\eta^\beta).
\end{aligned}}
\label{eq:matrix-beta-field}
\end{equation}
Here each all-ones outer product has shape
$C_{\mathrm{pr}}\times N_0$; hence every entry of $\balpha$ and
$\bbeta$ exceeds one.

At the mask-feature resolution, training draws the elementwise
reparameterized matrix sample
\begin{equation}
    \bm\Gamma_b\sim\operatorname{Beta}(\balpha,\bbeta),
    \qquad b=1,\ldots,B,
\end{equation}
independently across blocks, prompts, and mask-grid sites.  The same sampled
realization $\bm\Gamma_b$ is shared by all $Z$ slices of block $b$.  Validation
and deployment instead use
\begin{equation}
    \bar{\bm\Gamma}=\balpha\oslash(\balpha+\bbeta),
    \label{eq:matrix-beta-mean}
\end{equation}
where $\oslash$ is elementwise division.  The RBF construction makes the
distribution-parameter field smooth; it does not correlate the elementwise
training samples.  Area interpolation, rather than a second Beta evaluation,
produces $\bm\Gamma_b^r$ at every $r\in\cQ$ from this one mask-grid
realization.

Writing $\bm\gamma_{bc}^r\in\R^{N_r}$ for row $c$ of $\bm\Gamma_b^r$ and
$\bm F_{p,bz}^r,\bm F_{s,bz}^r\in\R^{N_r\times C_r}$, the visual matrix
presented to prompt $c$ is
\begin{equation}
\boxed{
    \widetilde{\bm F}_{bzc}^r
    =\bm F_{s,bz}^r
     +\diag(\bm\gamma_{bc}^r)
       (\bm F_{p,bz}^r-\bm F_{s,bz}^r).}
    \label{eq:matrix-spatial-fusion}
\end{equation}
This makes the broadcasting rule explicit: the gate varies with prompt and
position but is scalar across feature channels.  For example, a gate value
$0.8$ supplies $0.8\bm F_p+0.2\bm F_s$ at that site.

The prior has mean $m_0=0.7$ and concentration $\kappa_0=10$, i.e.
$(\alpha_0,\beta_0)=(7,3)$.  Both coefficient-head weights and head biases are
zero-initialized, while
\begin{equation}
 b_\alpha=\softplus^{-1}(7-1),\qquad
 b_\beta=\softplus^{-1}(3-1),
\end{equation}
so Eq.~\eqref{eq:matrix-beta-field} is exactly the homogeneous prior at
initialization.  Its matrix-averaged
regularizer is
\begin{equation}
    \cL_{\mathrm{route}}
    =\frac{1}{C_{\mathrm{pr}}N_{\mathrm{mask}}}
      \left\langle\ones_{C_{\mathrm{pr}}}\ones_{N_0}^\top,
      \KL\!\left[
       \operatorname{Beta}(\balpha,\bbeta)
       \,\|\,\operatorname{Beta}(7,3)
      \right]\right\rangle_F.
    \label{eq:matrix-route-kl}
\end{equation}
This average contains neither a batch nor a slice factor: the Beta parameter
field is defined once per prompt, and the block dimension enters only through
training-time samples.

All $BZC_{\mathrm{pr}}$ fused visual--prompt pairs are flattened in $(b,z,c)$
order and evaluated in one frozen BiomedParse predictor call.  If its query-logit tensor is
$\bm Q_{bzc}\in\R^{Q_m\times N_o}$, the implemented output is
\begin{equation}
    \bm L_{bzc}
    =\frac{1}{Q_m}\ones^\top\bm Q_{bzc},
    \label{eq:matrix-query-mean}
\end{equation}
followed by trilinear logit resizing to $(Z,H,W)$.  The vector consumed by the
segmentation functional below is
\begin{equation}
 \bm L_{bc}=\operatorname{vec}_{z,\bm u}\!\left(
   \Resize_{\mathrm{tri}}\{\bm L_{bzc}\}_{z=1}^{Z}
   \right)\in\R^{V},\qquad V=ZHW.
\end{equation}
Object-existence scores do
not weight this mean.  Importantly, the router is conditioned on frozen prompt
semantics and fixed coordinates, not on image features: for a given prompt,
different images share the same deterministic mean field at deployment.

\section{A Vectorized Prompt Segmentation Functional}

Let $\bm L_{bc},\bm y_{bc}\in\R^V$ collect the output logit and binary target
for prompt $c$ in block $b$.  For generic training,
$\bm y_{bc}=\operatorname{vec}(\ind[\bm Y_b=k_c])$; LGE uses the corresponding
column of the grouped target $\bm Y^j$.  The mask
$\bm m_b\in\{0,1\}^V$ removes both target $-1$ and slices with $v_{bz}=0$.
Define two masked scalar functionals
\begin{equation}
\begin{aligned}
 \mathfrak B(\bm L,\bm y;\bm m)
 &=\frac{\ones^\top\!\left[\bm m\odot
        (\softplus(\bm L)-\bm y\odot\bm L)\right]}
        {\max(\ones^\top\bm m,1)},\\
 \mathfrak D(\bm L,\bm y;\bm m)
 &=1-\frac{2(\bm m\odot\sigm(\bm L))^\top\bm y+10^{-6}}
 {\ones^\top(\bm m\odot\sigm(\bm L))
  +\ones^\top(\bm m\odot\bm y)+10^{-6}}.
\end{aligned}
\label{eq:vector-bce-dice}
\end{equation}
The first line is BCE-with-logits in closed form.  Put the empty-target and
prompt-validity matrices in
$\bm\Delta^{\mathrm{emp}},\bm\Omega
\in\{0,1\}^{B\times C_{\mathrm{pr}}}$,
where
\begin{equation}
 \Delta^{\mathrm{emp}}_{bc}=\ind\!\left[
 \bm m_b^\top\bm y_{bc}<5\times10^{-4}\,\ones^\top\bm m_b\right],
 \qquad
 W_{bc}=\Omega_{bc}\left(1-\tfrac12\Delta^{\mathrm{emp}}_{bc}\right).
\end{equation}
Define the pair-loss matrix $\bm S\in\R^{B\times C_{\mathrm{pr}}}$ by
\begin{equation}
    S_{bc}=\mathfrak B(\bm L_{bc},\bm y_{bc};\bm m_b)
       +(1-\Delta^{\mathrm{emp}}_{bc})
        \mathfrak D(\bm L_{bc},\bm y_{bc};\bm m_b).
    \label{eq:pair-loss-matrix}
\end{equation}
The generic and LGE reductions are now two matrix contractions of the same
$\bm S$:
\begin{equation}
\boxed{
\begin{aligned}
 \cL_{\mathrm{seg}}^{\mathrm{mean}}
 &=\frac{\langle\bm W,\bm S\rangle_F}
        {\max(\langle\bm W,
          \ones_B\ones_{C_{\mathrm{pr}}}^\top\rangle_F,10^{-6})},\\
 \cL_{\mathrm{seg}}^{\mathrm{LGE}}
 &=\ones_{C_{\mathrm{pr}}}^\top
   \left[(\bm W\odot\bm S)^\top\ones_B
   \oslash\max(\bm W^\top\ones_B,10^{-6})\right].
\end{aligned}}
\label{eq:matrix-seg-reductions}
\end{equation}
$\oslash$ and the vector-valued maximum in the second line act elementwise.
Thus generic training averages all valid block--prompt pairs, whereas LGE first
averages each prompt over blocks and then sums prompts.

\section{Two Branches, One Transport Variational Principle}

OT is solved independently for each valid slice and level.  After adaptive
average pooling to
$H_\ell^{\mathrm{OT}}=\min(H_\ell,32)$ and
$W_\ell^{\mathrm{OT}}=\min(W_\ell,32)$, let
\begin{equation}
 \bm X_{\ell m}^\xi,\bm Y_{\ell m}^\xi
 \in\R^{N_\ell^{\mathrm{OT}}\times C_\ell},
 \qquad \xi\in\{p,s\},\quad m=1,\ldots,M,
 \label{eq:ot-token-matrices}
\end{equation}
be foundation and channel-aligned expert token matrices.  The implementation
batches these independent problems but never mixes tokens from different
slices into one plan.  Because the expert feature is resized before adaptation,
both sides share the same pooled grid; pooling only downsamples and never
upsamples.  We suppress $(\ell,m)$ in the mass, cost, and plan definitions
below and restore both indices in the loss aggregation.

\subsection{Mass vectors as branch-specific matrix functionals}

All mass-stabilization operations below use $\epsnum=10^{-8}$, distinct from
the Sinkhorn temperatures $\eps_p$ and $\eps_s$.

For a valid slice, let $\bm M_c=\ind[\bm Y=k_c]$ and collect its pooled class
occupancy, boundary, and rescue maps as
the columns of $\bm O,\bm B,\bm M_{\mathrm{res}}
\in\R^{N_\ell^{\mathrm{OT}}\times C_{\mathrm{pr}}}$.  Occupancy
and boundary use adaptive average pooling; rescue uses adaptive max pooling.
With inverse-square-root present-class budget $\bm\pi$, the complete structural
score is the matrix--vector product
\begin{equation}
\begin{aligned}
 \bar\pi_c&=\begin{cases}
   |\bm M_c|^{-1/2},&|\bm M_c|>0,\\
   0,&|\bm M_c|=0,
  \end{cases}
 &\bm\pi&=\Norm(\bar{\bm\pi}),\\
 \bm h&=(\bm O+1.5\bm B+0.2\bm M_{\mathrm{res}})\bm\pi.
\end{aligned}
    \label{eq:matrix-structure-score}
\end{equation}
The boundary column is obtained from the clipped difference between a
$3\times3$ dilation and erosion.  The budget is normalized to sum to one; if
all classes are absent, the implementation uses its uniform fallback while the
zero structure score triggers the energy fallback below.

For a row-token matrix $\bm F$, define the stabilized magnitude and
mean-normalized energy vectors
\begin{equation}
    \bm\nu(\bm F)_i=\sqrt{\|\bm F_{i:}\|_2^2+\epsnum},
    \qquad
    \mathsf E(\bm F)
      =\bm\nu(\bm F)\big/\max(\operatorname{mean}\bm\nu(\bm F),\epsnum).
\end{equation}
Then the two balanced-$P$ marginal candidates are compactly
\begin{equation}
\begin{aligned}
 \bm q_b^p&=\bm h\odot[\ones+0.1\mathsf E(\bm X^p)],&
 \bm a^p&=\Norm(\bm q_b^p),\\
 \bm q_n^p&=\bm h\odot[\ones+0.1\mathsf E(\bm Y^p)],&
 \bm b^p&=\Norm(\bm q_n^p).
\end{aligned}
\label{eq:matrix-p-masses}
\end{equation}
If $\ones^\top\bm h\leq\epsnum$, $\bm q_b^p$ and $\bm q_n^p$ are replaced by
their respective energy vectors.  Every $\Norm$ divides by total mass and uses
a uniform vector when that total is at most $\epsnum$.

For the residual branch, define the specificity functional
\begin{equation}
    \bm\chi(\bm P,\bm S)
    =\bm\nu(\bm S)\oslash
      [\bm\nu(\bm P)+\bm\nu(\bm S)+\epsnum\ones].
\end{equation}
Let $\bm e_c^b$ and $\bm e_c^n$ be the detached pixel-wise BCE maps of the
foundation and expert for prompt $c$.  Before pooling, both sides use the same
prompt-validity reduction
\begin{equation}
 \bm d_{\mathrm{pix}}
 =\frac{\sum_c\Omega_c\bm e_c^b}{\max(\sum_c\Omega_c,1)},
 \qquad
 \bm e^n_{\mathrm{pix}}
 =\frac{\sum_c\Omega_c\bm e_c^n}{\max(\sum_c\Omega_c,1)}.
 \label{eq:valid-prompt-error-maps}
\end{equation}
Adaptive average pooling gives foundation difficulty $\bm d$ and expert error
$\bm e^n$ on the OT grid.  With advantage
$\bm\Delta_{\mathrm{adv}}=\bm d-\bm e^n$, the code's dense smooth scores are
\begin{equation}
\begin{aligned}
 \widehat{\bm d}
 &=\clip\!\left(
    \bm d/\max(\operatorname{mean}\bm d,\epsnum),0,3\right),\\
 s_{\Delta}&=\max(\operatorname{mean}|\bm\Delta_{\mathrm{adv}}|,\epsnum),\qquad
 \bm\zeta=\clip\!\left(
  2\sigm(\bm\Delta_{\mathrm{adv}}/(0.5s_{\Delta})),0,3\right).
\end{aligned}
\label{eq:matrix-s-gain}
\end{equation}
Hence $\zeta_i>1$ favors sites at which the expert has lower error, whereas
$\zeta_i<1$ suppresses sites at which it is worse.  The unbalanced-$S$ marginals
are
\begin{equation}
\begin{aligned}
 \bm a^s&=\Norm\!\left(
   \mathsf E(\bm X^s)\odot\bm\chi(\bm X^p,\bm X^s)
   \odot(0.1\ones+\widehat{\bm d})\right),\\
 \bm b^s&=\Norm\!\left(
   \mathsf E(\bm Y^s)\odot\bm\chi(\bm Y^p,\bm Y^s)
   \odot(0.1\ones+\bm\zeta)\right).
\end{aligned}
\label{eq:matrix-s-masses}
\end{equation}
For grouped LGE labels, the expert binary logit used in these error maps is
$\operatorname{logit}(\clip(\sum_{k\in G_c}\operatorname{softmax}
(\bm L^n)_k,10^{-6},1-10^{-6}))$.

\subsection{Cost matrices}

Let $\overline{\bm X}$ and $\overline{\bm Y}$ denote detached row-wise
$\ell_2$ normalization, using denominator $\max(\|\bm x_i\|_2,10^{-6})$ so
that a zero row remains zero.  If $\bm U_b,\bm U_n$ collect normalized spatial
coordinates, define
\begin{equation}
\begin{aligned}
 \bm C_{\mathrm{feat}}^\xi
   &=\relu{\ones_{N_\ell^{\mathrm{OT}}}
             \ones_{N_\ell^{\mathrm{OT}}}^\top-
        \overline{\bm X}^{\xi}\overline{\bm Y}^{\xi\top}},\\
 \bm C_{\mathrm{coord}}
   &=\relu{\operatorname{Dist}(\bm U_b,\bm U_n)-0.25}^{\odot2},\\
 \bm C_{\mathrm{sem}}
   &=\relu{\ones_{N_\ell^{\mathrm{OT}}}
             \ones_{N_\ell^{\mathrm{OT}}}^\top-
        \overline{\bm M}_b\overline{\bm M}_n^\top}.
\end{aligned}
\label{eq:matrix-transport-costs}
\end{equation}
The label rows in $\overline{\bm M}$ are stabilized $\ell_1$ normalizations
over classes.  The two implemented costs are the single branch-indexed family
\begin{equation}
    \boxed{
    \bm C^\xi
    =\bm C_{\mathrm{feat}}^\xi
      +0.25\bm C_{\mathrm{coord}}
      +0.25\ind[\xi=p]\bm C_{\mathrm{sem}},
    \qquad \xi\in\{p,s\}.}
    \label{eq:unified-branch-cost}
\end{equation}

\subsection{One branch-indexed transport problem}

Define $\Ent(\bm\Pi)=\sum_{ij}\Pi_{ij}(\log\Pi_{ij}-1)$ and the branch
regularizers
\begin{equation}
\begin{aligned}
 \mathfrak P_p(\bm\Pi;\bm a^p,\bm b^p)
 &=\begin{cases}
   0,&\bm\Pi\ones=\bm a^p,\ \bm\Pi^\top\ones=\bm b^p,\\
   +\infty,&\text{otherwise},
   \end{cases}\\
 \mathfrak P_s(\bm\Pi;\bm a^s,\bm b^s)
 &=\rho_b\KL(\bm\Pi\ones\,\|\,\bm a^s)
   +\rho_n\KL(\bm\Pi^\top\ones\,\|\,\bm b^s).
\end{aligned}
\end{equation}
The two KL terms are generalized KL divergences and therefore accept
non-unit-mass UOT marginals.
Balanced OT and UOT are therefore two instances of one variational equation:
\begin{equation}
\boxed{
    \bm\Pi^\xi
    =\arg\min_{\bm\Pi\geq0}
      \left\{
      \langle\bm\Pi,\bm C^\xi\rangle_F
      +\eps_\xi\Ent(\bm\Pi)
      +\mathfrak P_\xi(\bm\Pi;\bm a^\xi,\bm b^\xi)
      \right\},
    \quad \xi\in\{p,s\}.}
    \label{eq:master-transport}
\end{equation}
The defaults are $\eps_p=\eps_s=0.1$, $\rho_b=1.0$, and $\rho_n=0.2$.
Both plans use 30 float32 log-space Sinkhorn iterations; for $S$, the relaxation powers
are $\tau_b=\rho_b/(\rho_b+\eps_s)$ and
$\tau_n=\rho_n/(\rho_n+\eps_s)$.

Let $\bm\varrho^\xi=\bm\Pi^\xi\ones$ and
$\bm D_\xi=\diag(\max(\bm\varrho^\xi,10^{-8}))$.  The expert-to-foundation
barycentric teacher is the matrix projection
\begin{equation}
    \widetilde{\bm Y}^{\xi}
    =\bm D_\xi^{-1}\bm\Pi^\xi\bm Y^\xi.
    \label{eq:matrix-barycentric-teacher}
\end{equation}
The plan is built from detached normalized copies, whereas $\bm Y^\xi$ in
Eq.~\eqref{eq:matrix-barycentric-teacher} is the raw pooled expert feature.
For the row-wise cosine discrepancy vector
$\bm\delta(\bm X,\bm Y)=\ones-\rowcos(\bm X,\bm Y)$, the two distillation
terms share one formula,
\begin{equation}
\boxed{
 \cL_{\xi,\ell}
 =\frac{\displaystyle\sum_{m=1}^{M}
   \bm w_{\ell m}^{\xi\top}
   \bm\delta\!\left(\bm X_{\ell m}^\xi,
      \stopgrad(\widetilde{\bm Y}_{\ell m}^{\xi})\right)}
   {\displaystyle\max\!\left(
      \sum_{m=1}^{M}\ones^\top\bm w_{\ell m}^\xi,10^{-8}\right)},
 \quad
 \bm w_{\ell m}^p=\bm a_{\ell m}^p,\quad
 \bm w_{\ell m}^s=\bm\Pi_{\ell m}^s\ones.}
 \label{eq:unified-transport-distillation}
\end{equation}
The final branch losses are
\begin{equation}
 \cL_{P\text{-OT}}=\frac14\sum_{\ell\in\cI}\cL_{p,\ell},
 \qquad
 \cL_{S\text{-UOT}}=\frac14\sum_{\ell\in\cI}\cL_{s,\ell}.
\end{equation}
Because every $\bm a^p_{\ell m}$ sums to one, Eq.~\eqref{eq:unified-transport-distillation}
gives valid slices equal weight in $P$; $S$ slices are instead weighted by
their received UOT mass.  If the layer-wide received total
$\sum_m\ones^\top\bm w^s_{\ell m}$ is at most $10^{-6}$, that layer's $S$
loss is zero.  Masses, error maps, costs, plans, and barycentric
teachers are detached.  Consequently, OT gradients update only the raw
foundation-side factor matrices; expert adapters are trained by reconstruction
and separation, not by differentiation through Sinkhorn.

\section{One Master Objective}

Collect the six executable losses into
\begin{equation}
    \bm{\lambda}(e)=
    \begin{bmatrix}
      \lambda_{\mathrm{seg}}&\lambda_{\mathrm{ae}}&
      \lambda_{\mathrm{ort}}&\lambda_{\mathrm{route}}(e)&
      \lambda_p(e)&\lambda_s(e)
    \end{bmatrix}^{\!\top},
    \qquad
    \bm{\mathcal J}_\Theta=
    \begin{bmatrix}
      \cL_{\mathrm{seg}}&\cL_{\mathrm{ae}}&
      \cL_{\mathrm{ort}}&\cL_{\mathrm{route}}&
      \cL_{P\text{-OT}}&\cL_{S\text{-UOT}}
    \end{bmatrix}^{\!\top}.
\end{equation}
The complete generic objective is simply
\begin{equation}
\boxed{
\begin{aligned}
 \cL(\Theta;e)
 &=\bm{\lambda}(e)^\top\bm{\mathcal J}_\Theta\\
 &=\underbrace{
   \lambda_{\mathrm{seg}}\cL_{\mathrm{seg}}
   +\lambda_{\mathrm{route}}\cL_{\mathrm{route}}
   }_{\text{prompt-conditioned prediction}}
  +\underbrace{
   \lambda_{\mathrm{ae}}\cL_{\mathrm{ae}}
   +\lambda_{\mathrm{ort}}\cL_{\mathrm{ort}}
   }_{\text{information-preserving complementary factors}}\\
 &\quad
  +\underbrace{
   \lambda_p\cL_{P\text{-OT}}
   +\lambda_s\cL_{S\text{-UOT}}
   }_{\text{component-specific expert transfer}}.
\end{aligned}}
    \label{eq:master-objective}
\end{equation}
Its target weights are
$(3,0.2,0.3,10^{-3},0.1,0.1)$.  Define the delayed ramp
\begin{equation}
    \operatorname{ramp}(e;e_0,E_w)
    =\begin{cases}
      0,&e<e_0,\\
      \min\!\left(1,\dfrac{e-e_0+1}{E_w}\right),&e\geq e_0.
     \end{cases}
\end{equation}
Then the exact schedules are
\begin{equation}
\begin{aligned}
 \lambda_{\mathrm{route}}(e)&=10^{-3}\min(1,e/5),\\
 \lambda_p(e)&=0.1\operatorname{ramp}(e;2,5),\qquad
 \lambda_s(e)=0.1\operatorname{ramp}(e;3,5).
\end{aligned}
\label{eq:unified-schedules}
\end{equation}
AdamW jointly updates all of $\Theta$ with learning rate $10^{-4}$, weight
decay $10^{-4}$, optional automatic mixed precision, and global gradient-norm
clipping at $5$.  Generic training keeps a constant learning rate.  An
$E$-epoch LGE run uses
\begin{equation}
 \eta_e=\eta_0\left[0.05+0.95
 \frac{1+\cos(\pi(e-1)/(E-1))}{2}\right].
\end{equation}

\section{LGE as a Composition of the Same Student Operator}

The LGE pathway does not introduce a second network.  Let
$\mathfrak S_\Theta(\bm x^b;\cG)=\bm L$ denote the pure-student prediction map
in Eq.~\eqref{eq:master-student} for prompt groups $\cG$.  LGE composes the
same $\mathfrak S_\Theta$ once globally and once locally.  BiomedParse inputs reuse
the frozen nnUNet crop and resampling geometry, but replace intensity
normalization by
\begin{equation}
 x_i^b=255\,
 \frac{\clip(x_i,q_1,q_{99})-q_1}{q_{99}-q_1},
 \label{eq:lge-normalization}
\end{equation}
where $q_1,q_{99}$ are the first and $99$th percentiles after removing values at
or below the volume's $0.1$th percentile.  Degenerate ranges use
$q_{99}=q_1+1$; output is clipped to $[0,255]$ and stored as uint8.

For source labels scar $1$, edema $2$, LV $3$, normal myocardium $4$, and RV
$5$, the prompt-group tuples are
\begin{equation}
\begin{aligned}
 \cG^A&=((3),(5),(1,2,4)),\\
 \cG^R_{\mathrm{split}}&=((3),(5),(4),(1),(2)),\\
 \cG^R_{\mathrm{merged}}&=((3),(5),(4),(1,2)).
\end{aligned}
\label{eq:lge-groups}
\end{equation}
Thus $C_A=3$, while $C_R=5$ for V3 split pathology and $C_R=4$ for V2 merged
pathology.  Group remapping is performed independently for each pass, and the
source-label groups within a pass are required to be disjoint.  Equivalently,
if $\bm Y_{\mathrm{src}}\in\{0,1\}^{V\times5}$ is the raw one-hot target and
$[\bm G_j]_{kc}=\ind[k\in G_c^j]$, then the complete remapping is the matrix
product
\begin{equation}
    \bm Y^j=\bm Y_{\mathrm{src}}\bm G_j,
    \qquad \bm G_j\in\{0,1\}^{5\times C_j},\quad j\in\{A,R\}.
\end{equation}
Background voxels are all-zero rows in $\bm Y_{\mathrm{src}}$; target $-1$
rows are removed by the valid mask rather than assigned to a foreground group.

For validation and deployment, the two predictions form the online composition
\begin{equation}
\boxed{
\begin{aligned}
 \bm L^A
   &=\mathfrak S_\Theta(\bm x^b;\cG^A),\\
 \mathcal R^A_{\mathrm{online}}
   &=\stopgrad\!\left\{\RoiBox_{\gamma=1.25}\!\left(
      \ind[\sigm(\bm L^A_{\mathrm{MYO}})\geq0.3]
      \right)\right\},\\
 \bm L^R
   &=\mathfrak S_\Theta(
       \Crop(\bm x^b;\mathcal R^A_{\mathrm{online}});\cG^R).
\end{aligned}}
\label{eq:two-pass-composition}
\end{equation}
The box is the smallest axis-aligned foreground box, enlarged about its center,
clipped/shifted to the image boundary, and constrained to a minimum side of
$8$ pixels.  An empty mask falls back to the full image.  Cropped images are
bilinearly resized to $256^2$ and labels by nearest-neighbor interpolation.
During Pass~R training, exactly the same ROI coordinates are applied to the
nnUNet input, BiomedParse input, and target before their respective
interpolation rules; pure-student deployment has no nnUNet input.

Training uses Pass~A alone for the first ten epochs.  Before epoch $11$,
deterministic student predictions create a detached ROI cache; the default
refresh interval is zero.  Validation and deployment instead construct the ROI
online.  To distinguish this training-only computation from pure prediction,
define
\begin{equation}
 \mathfrak J_\Theta(\bm x^b,\bm x^n,\bm Y,\bm V;\cG)
 =\bigl(\bm L,\cL_{\mathrm{seg}},\cL_{\mathrm{aux}}\bigr),
 \label{eq:training-evaluator}
\end{equation}
where $\mathfrak J_\Theta$ evaluates the pure-student logits together with the
expert-dependent training losses.  Pass~A applies it to the full aligned
tensors.  Mixed-training Pass~R instead uses
\begin{equation}
 \mathcal R^{\mathrm{train}}_b
   =\operatorname{Jitter}(\mathcal R^{\mathrm{cache}}_b),
 \qquad
 \mathfrak J_\Theta\!\left(
   \Crop_{\mathcal R^{\mathrm{train}}_b}
   (\bm x^b,\bm x^n,\bm Y,\bm V);\cG^R\right).
 \label{eq:cached-training-roi}
\end{equation}
Thus the current Pass~A logit does not select the mixed-training crop.  If group
$G_c$ is present in the full target, its refinement prompt is
valid only when
\begin{equation}
    \Omega^R_{bc}
    =\ind\!\left[
      \frac{\|\ind[Y\in G_c]\odot\ind[\mathcal R^{\mathrm{used}}_b]\|_1}
           {\|\ind[Y\in G_c]\|_1}\geq0.01\right];
\end{equation}
Here $\mathcal R^{\mathrm{used}}_b$ is the cached/jittered ROI in training and
the online ROI otherwise; an absent full-slice group remains a valid negative.

The warm-up and mixed stages share the auxiliary functional below.  Their
epoch-dependent objective is
\begin{equation}
 \boxed{\cL_{\mathrm{LGE}}(e)=\begin{cases}
   \cL^A_{\mathrm{seg}}+\cL^A_{\mathrm{aux}},&e\leq10,\\[1mm]
   \cL^A_{\mathrm{seg}}+\cL^R_{\mathrm{seg}}
   +\dfrac12(\cL^A_{\mathrm{aux}}+\cL^R_{\mathrm{aux}}),&e\geq11.
 \end{cases}}
\label{eq:unified-lge-objective}
\end{equation}
where, for $j\in\{A,R\}$,
\begin{equation}
 \cL^j_{\mathrm{aux}}
 =0.2\cL^j_{\mathrm{ae}}+0.3\cL^j_{\mathrm{ort}}
  +\lambda_{\mathrm{route}}(e)\cL^j_{\mathrm{route}}
  +\lambda_p(e)\cL^j_{P\text{-OT}}
  +\lambda_s(e)\cL^j_{S\text{-UOT}}.
\end{equation}
The implemented pass weights are $\lambda_A=\lambda_R=1$.  Pass~R is not
evaluated in the first case.  In the second, both segmentation losses remain
unscaled, both auxiliary groups receive factor $1/2$, and both forwards
accumulate gradients before one optimizer update.

During mixed training, the cached ROI may shift by $3\%$ of its extent and
enlarge by a factor in $[1,1.15]$.  The two passes then receive independent
affine/intensity transforms: rotation $\pm10^\circ$, scale $[0.95,1.05]$,
translation $\pm5\%$, horizontal/vertical flip probabilities $0.5/0.2$, and
an intensity perturbation with probability $0.8$.  Contrast is in $[0.8,1.2]$,
shift in $[-0.1,0.1]$ times the slice standard deviation, Gaussian-noise scale
in $[0,0.04]$ times that deviation, and the BiomedParse gamma in
$[0.75,1.35]$.

\subsection{Hard switching as a matrix selector}

Let $\bm\iota_R\in\{0,1\}^{V}$ be the rectangular ROI indicator and
$\bm D_R=\diag(\bm\iota_R)$.  Let $\bm L^R_{\uparrow}\in\R^{V\times C_R}$ denote
the refinement logits after bilinear resizing to the ROI extent and insertion
into the full model lattice.  Define the fixed channel-lifting matrix
$\bm\Lambda_{AR}\in\{0,1\}^{3\times C_R}$ by
$[\bm\Lambda_{AR}]_{11}=[\bm\Lambda_{AR}]_{22}=1$ and zero elsewhere, and let
$\bm q_{AR}=(1,1,0,\ldots,0)^\top$.  The outside-ROI fallback is
\begin{equation}
 \bm L^{\mathrm{out}}
 =\bm L^A\bm\Lambda_{AR}
   -20\ones_V(\ones_{C_R}-\bm q_{AR})^\top
 \in\R^{V\times C_R}.
\end{equation}
The current V2/V3 rule is the matrix selector
\begin{equation}
    \boxed{
    \bm L^{\mathrm{final}}
    =\bm D_R\bm L^R_{\uparrow}
     +(\bm I-\bm D_R)\bm L^{\mathrm{out}}.}
    \label{eq:matrix-hard-switch}
\end{equation}
Thus all refinement classes are used inside the ROI; outside it, Pass~A supplies
LV/RV and every myocardial/pathology foreground logit is clamped to $-20$.
There is no averaging or soft blending.  After prepending a zero background
column, the composite logit tensor is bilinearly restored to aligned geometry
and then resampled with the nnUNet probability-resampling function.  Class
argmax is taken only in that pre-cropping geometry, after which the discrete
segmentation is inserted into the saved crop box and inverse-transposed to raw
NIfTI geometry.

\section{Training--Deployment Boundary and Current Instantiation}

At deployment, the retained parameter set is
\begin{equation}
    \Theta_{\mathrm{deploy}}
    =\left\{\bm W^b_{\ell p},\bm W^b_{\ell s}
      \right\}_{\ell\in\cI}\cup\Theta_{\mathrm{router}}.
\end{equation}
The nnUNet expert, expert autoencoders, expert preprocessing tensors, mass and
cost builders, and Sinkhorn solvers are absent.  BiomedParse features are
factorized by Eq.~\eqref{eq:foundation-block-factorization}, fused with the
deterministic mean in Eq.~\eqref{eq:matrix-beta-mean}, and sent to the frozen
predictor.  LGE deployment invokes this pure-student map twice through
Eq.~\eqref{eq:two-pass-composition}.

For the current $256\times256$ LGE instantiation, the spatial feature sizes are
$(64,32,16,8)^2$ at levels $2$--$5$.  The corresponding nnUNet channels are
$(128,256,512,512)$ and BiomedParse channels are
$(192,384,768,1536)$, yielding OT grids $(32,32,16,8)^2$.  The complete
training modules contain $10{,}140{,}994$ trainable parameters; the deployment
subset contains $6{,}432{,}130$.

\end{document}
